%% ===================================================================
\section{\AEGIS Construction}
\label{sec:framework}
%% ===================================================================

Translating the results of \cref{sec:theory} into a practical tool requires overcoming two computational barriers: the full sensitivity matrix $S \in \R^{Nd \times N^2}$ is too large to materialize, and the resolvent $(I - J_z)^{-1}$ is too expensive to form directly. The \AEGIS pipeline addresses both through matrix-free computation.

\subsection{Architecture Overview}

\AEGIS takes as input a trained GNN (any architecture whose message passing is differentiable with respect to edge weights) together with a target graph and produces three outputs: (i)~a per-edge vulnerability spectrum ranking edges by adversarial impact, (ii)~the maximally sensitive perturbation direction via SVD, and (iii)~per-node first-order sensitivity thresholds. For contractive implicit GNNs (IGNN-class satisfying A1--A3), the pipeline additionally computes the critical budget $\ecrit$ and convergence diagnostics; these formal guarantees are specific to the implicit case. The pipeline operates in two regimes:
\begin{itemize}
\item \textbf{Dense path} ($N \leq 200$): materializes $S_c \in \R^{Nd \times |E|}$, computes exact SVD and unconstrained $r_v$. Applicable to subgraph-level analysis and small power grids.
\item \textbf{Matrix-free path} ($N > 200$, validated up to $N{=}7{,}650$): operates via JVP/VJP queries on $S_c$ without materializing any matrix. Computes randomized SVD, per-edge column norms, and constrained $r_v^{(c)}$ (\cref{prop:radius}).
\end{itemize}
The four stages are illustrated in \cref{fig:pipeline}.

\subsection{Stage 1: Subgraph Extraction and Forward Pass}

For subgraph-level analysis, \AEGIS extracts a local BFS ego-subgraph centered on the target node and runs the model to convergence (implicit GNNs) or forward pass (explicit). Localization is justified empirically (\cref{sec:subgraph_ablation}). The matrix-free pipeline also supports full-graph analysis without subgraph extraction (\cref{sec:scalability}).

\subsection{Stage 2: Matrix-Free Sensitivity Computation}

\AEGIS requires the action of the constrained sensitivity operator $S_c$ on vectors, without materializing the full matrix $S \in \R^{Nd \times N^2}$. For an edge perturbation vector $v \in \R^{|E|}$, the operator computes $S_c v = (I - J_z)^{-1} J_A P_c v$, where $P_c$ maps edge weights to symmetric adjacency perturbations, $J_A$ is the structural Jacobian, and $J_z = \partial F / \partial Z$ is the state Jacobian. Two ingredients make this matrix-free:

\textbf{Neumann-series resolvent.} The resolvent $(I - J_z)^{-1} b$ is evaluated via the truncated Neumann series $\sum_{k=0}^{K} J_z^k b$, which converges geometrically when $\kappa < 1$ (Assumption~A3). Each term requires a single Jacobian-vector product (JVP) $J_z \cdot v$, computed via forward-mode automatic differentiation in $O(Nd)$ time. The depth $K$ is set adaptively from a power-iteration estimate of $\kappa$, with early stopping when $\norm{J_z^k b} < 10^{-6} \norm{b}$.

\textbf{Autograd structural JVPs.} The adjacency Jacobian $J_A \cdot \mathrm{vec}(\delta A)$ is computed as a single forward-mode JVP of $F$ with respect to $A$ in direction $\delta A$, avoiding the $O(N^2)$ finite-difference columns required by the dense path.

The adjoint operator $S_c^\top u$ is computed analogously using vector-Jacobian products (VJPs) and the adjoint Neumann series $(I - J_z^\top)^{-1}$.

For small subgraphs ($N \leq 200$), the dense path (forming $J_z$, $J_A$ explicitly and solving via \texttt{torch.linalg.solve}) remains available and gives exact results.

\subsection{Stage 3: Constrained Projection}

The constraint projection $P_c$ (\cref{sec:theory}) maps edge vectors to symmetric adjacency perturbations ($\delta A_{ij} = \delta A_{ji} = v_k$). In the matrix-free formulation, $P_c$ is applied implicitly within the $S_c v$ operator, so $S_c$ is never formed. This $N^2 \to |E|$ reduction is what enables tight first-order predictions (\cref{sec:cross_domain,sec:explicit_extension}).

\subsection{Stage 4: Vulnerability Analysis}

From $S_c$, \AEGIS extracts three complementary outputs:

\textbf{Maximally sensitive perturbation direction.} A randomized SVD~\cite{halko2011finding} of the $S_c$ operator (using only matrix-vector products, with rank $k = 10$, oversampling $p = 10$, and $n_{\text{iter}} = 2$ power iterations) yields the leading singular triplet $(u_1, \sigma_1, v_1)$. The first-order maximally sensitive perturbation is $\delta A^* = \varepsilon \cdot v_1$, achieving a predicted equilibrium shift of $\varepsilon \cdot \sigma_1(S_c)$. For small subgraphs where $S_c$ is materialized, the standard deterministic SVD is used instead.

\textbf{Per-edge vulnerability spectrum.} For each edge $(i,j) \in E$, the vulnerability score $v_{ij} = \norm{[S_c]_{:,k}}_2$ measures how much a unit perturbation of that edge shifts the equilibrium. Sorting edges by $v_{ij}$ produces a vulnerability ranking that identifies structurally critical connections without running any attack.

\textbf{Per-node first-order sensitivity thresholds.} For node $v$ with classification margin $m_v$ (gap between the predicted and runner-up class logits), the first-order sensitivity threshold from \cref{prop:radius} is computed as follows. In the \emph{dense path} ($N \leq 200$), we use the unconstrained $r_v = m_v / (\norm{W_{y_v} - W_{c^*}}_2 \cdot \norm{S_v}_2)$, where $S_v$ denotes the block-rows of the materialized $S$. In the \emph{matrix-free path} ($N > 200$), the full $S$ is never formed; we instead compute the tighter constrained threshold $r_v^{(c)} = m_v / (\norm{W_{y_v} - W_{c^*}}_2 \cdot \norm{S_{c,v}}_2)$, where $\norm{S_{c,v}}_2$ is estimated via 10 iterations of randomized power iteration on the $S_c$ operator restricted to node~$v$'s block-rows. Since $\norm{S_{c,v}}_2 \leq \norm{S_v}_2$, the matrix-free thresholds are tighter than the dense-path ones. Any perturbation with $\norm{\delta \Ahat}_F < r_v$ (or $r_v^{(c)}$) preserves the classification of node $v$ to first order. These are diagnostic sensitivity thresholds, not probabilistic certificates (\cref{rem:certificates}).

\subsection{Computational Cost}

\textbf{Matrix-free pipeline.} Each $S_c v$ or $S_c^\top u$ evaluation costs $O(K \cdot Nd)$, where $K$ is the Neumann depth (set adaptively from $\kappa$; typically $K = 20$--$50$ for $\kappa < 0.8$, up to $\sim\!340$ for $\kappa \approx 0.96$). The randomized SVD performs $O((k + p) \cdot n_{\text{iter}})$ such evaluations, where $k$ is the number of requested singular values, $p$ the oversampling, and $n_{\text{iter}}$ the power iterations. Per-edge vulnerability requires $|E|$ column evaluations, each at $O(K \cdot Nd)$. Memory is $O(Nd)$ per operation; no large matrix is stored.

\textbf{Scalability.} On a single NVIDIA RTX 4090 (24\,GB), the matrix-free pipeline completes in 1.1s for $N{=}50$, 4.5s for $N{=}200$, 24s for $N{=}1{,}000$, and 78s for full Cora ($N{=}2{,}708$, $Nd{=}173{,}312$) using 1\,GB GPU memory (\cref{sec:scalability}). The dense path runs faster for $N \leq 100$ (0.7s at $N{=}50$) but exceeds 24\,GB at $N{=}500$.

The matrix-free formulation removes the $O((Nd)^3)$ linear-solve bottleneck and the $O(N^3 d)$ storage requirement that previously limited analysis to $N \approx 300$, enabling full-graph vulnerability analysis on graphs with thousands of nodes (validated up to $N{=}7{,}650$; \cref{sec:scalability}).
